\scanpage{27}
\[
\mathord{\text{\begin{tikzpicture}[x=1ex,y=1ex,baseline=0,line width=.3pt]
\draw(0,0)--(.60,0)..controls(.45,.8)and(.48,1.5)..(1.0,1.5)..controls(1.65,1.5)and(1.70,.6)..(1.25,.02)--(1.8,.02);
\draw(0,-.05)--(.80,-.05)..controls(.65,.8)and(.65,1.25)..(1.0,1.25)..controls(1.45,1.25)and(1.45,.65)..(1.18,-.05)--(1.8,-.05);
\end{tikzpicture}}}=\uline{\operatorname{End}}(\mathcal E,\mathcal E)\subset\uline{\operatorname{Hom}}(\mathcal E,\mathrm{Ens})=\widehat{\mathcal E^o}
\]
$\mathord{\text{\begin{tikzpicture}[x=1ex,y=1ex,baseline=0,line width=.3pt]
\draw(0,0)--(.60,0)..controls(.45,.8)and(.48,1.5)..(1.0,1.5)..controls(1.65,1.5)and(1.70,.6)..(1.25,.02)--(1.8,.02);
\draw(0,-.05)--(.80,-.05)..controls(.65,.8)and(.65,1.25)..(1.0,1.25)..controls(1.45,1.25)and(1.45,.65)..(1.18,-.05)--(1.8,-.05);
\end{tikzpicture}}}$ est la sous-catégorie (pas pleine!) de $\widehat{\mathcal E^o}$ \unsure{formée}

a) comme objets,
des foncteurs $F:\mathcal E\mathbin{\text{---}}(\mathrm{Ens})$ qui se factorisent
par la sous-catégorie $\mathcal E$ de $\mathrm{Ens}$, i.e. tels que

$1^\circ$) $F(E)$ est fini pour tt $E\in\operatorname{Ob}\mathcal E$

$2^\circ$) $F(u)$ est injectif $\forall(u:E\mathbin{\text{---}}F)\in\operatorname{Fl}(\mathcal E)$

b) comme morphismes $F\xrightarrow{u}G$, des seuls
hom fonctoriels tels que $u(E):F(E)\mathbin{\text{---}}G(E)$ soit
injectif pour tt $E\in\operatorname{Ob}\mathcal E$.

$\mathord{\text{\begin{tikzpicture}[x=1ex,y=1ex,baseline=0,line width=.3pt]
\draw(0,0)--(.60,0)..controls(.45,.8)and(.48,1.5)..(1.0,1.5)..controls(1.65,1.5)and(1.70,.6)..(1.25,.02)--(1.8,.02);
\draw(0,-.05)--(.80,-.05)..controls(.65,.8)and(.65,1.25)..(1.0,1.25)..controls(1.45,1.25)and(1.45,.65)..(1.18,-.05)--(1.8,-.05);
\end{tikzpicture}}}$ a des opérations $\amalg,\times,\circ$
\[
\left\{\begin{aligned}
(F\unsure{+|\amalg}G)(E)&=F(E)\amalg G(E)\\
(F\times G)(E)&=F(E)\times G(E)\\
(F\circ G)(E)&=F(G(E))
\end{aligned}\right.
\]
\mspage{26}
